\section{Estratégias de crescimento}

\subsection{Estratégia 1: piloto territorial com operação assistida}

Na matriz de Ansoff, trata-se de desenvolvimento de mercado controlado [11]: testar o fluxo em Belo Horizonte, começando por um bairro ou rede comunitária que a equipe consiga acompanhar. Antes do piloto, a equipe deve entrevistar 15 a 20 tutores, conversar com cinco clínicas e mapear duas organizações comunitárias, seguindo o desenvolvimento de clientes [6]. Esses números são metas de pesquisa propostas, não contatos já realizados.

Em seguida, a equipe revisa a explicação dos créditos e do QR Code e testa tarefas com participantes operando o aplicativo sozinhos. Depois, executa por até oito semanas um piloto \textit{concierge}, em que cadastro de parceiros e encaminhamentos são acompanhados manualmente, prática recomendada para aprender antes de automatizar [16]. Um grupo de 40 a 50 tutores e duas ou três clínicas é hipótese de escala de teste, condicionada a consentimento e oferta real.

\subsection{Estratégia 2: aquisição local baseada em confiança e parcerias}

A segunda estratégia combina marketing de aquisição e parcerias: crescer por conteúdo útil e indicação de clínicas, ONGs e protetores, em vez de comprar alcance antes de provar ativação. Weinberg e Mares recomendam testar poucos canais em paralelo e concentrar esforço no que funciona [14]; Ellis e Brown propõem ciclos curtos de experimentos sobre o funil [15]. Um QR ou link por parceiro permite rastrear a origem do cadastro.

A métrica principal não é o número de downloads. Ativação significa criar o perfil do pet e registrar um lembrete em até sete dias; retenção D30 significa realizar uma ação relevante entre os dias 22 e 30. A recomendação (NPS) é medida ao fim do piloto [13]. Campanhas pagas só entram quando o custo por tutor ativado for conhecido.

\begin{table}[H]
\centering\small
\caption{Funil de crescimento proposto (metas de aprendizagem, não resultados).}
\begin{tabularx}{\textwidth}{P{3.2cm}YP{3.2cm}}
\toprule
\textbf{Etapa} & \textbf{Definição} & \textbf{Meta do piloto} \\
\midrule
Aquisição & Cadastro vindo de parceiro ou evento rastreado. & 40--50 tutores convidados \\
Ativação & Perfil do pet + primeiro lembrete em 7 dias. & 60\% dos cadastrados \\
Retenção & Ação relevante entre os dias 22 e 30. & 30\% dos ativados \\
Recomendação & NPS ao final do piloto. & Registrar e comparar \\
Lado da oferta & Clínicas que aceitam novo ciclo. & 2 clínicas \\
\bottomrule
\end{tabularx}
\end{table}
